\chapter{References and source inventory}
\label{ch:references}
\index{references}
\index{source inventory}

This manual cites the reviewed codebase directly and uses the following MPM literature as the initial formal reference set for Chapter~\ref{ch:theory}.  The list emphasizes the requested Homel-Brannon-Guilkey, Nairn, Bardenhagen, and Sulsky literature families.  Future versions should move these entries to BibTeX or BibLaTeX so the LaTeX report and later Sphinx documentation can share one reference database.

\begin{thebibliography}{99}

\bibitem{sulsky1994history}
D. Sulsky, Z. Chen, and H. L. Schreyer.
\newblock A particle method for history-dependent materials.
\newblock \emph{Computer Methods in Applied Mechanics and Engineering}, 118(1--2):179--196, 1994.
\newblock \url{https://doi.org/10.1016/0045-7825(94)90112-0}.

\bibitem{bardenhagen2004gimp}
S. G. Bardenhagen and E. M. Kober.
\newblock The generalized interpolation material point method.
\newblock \emph{Computer Modeling in Engineering \& Sciences}, 5(6):477--496, 2004.
\newblock \url{https://doi.org/10.3970/cmes.2004.005.477}.

\bibitem{sadeghirad2011cpdi}
A. Sadeghirad, R. M. Brannon, and J. Burghardt.
\newblock A convected particle domain interpolation technique to extend applicability of the material point method for problems involving massive deformations.
\newblock \emph{International Journal for Numerical Methods in Engineering}, 86(12):1435--1456, 2011.
\newblock \url{https://doi.org/10.1002/nme.3110}.

\bibitem{sadeghirad2013cpdi2}
A. Sadeghirad, R. M. Brannon, and J. E. Guilkey.
\newblock Second-order convected particle domain interpolation (CPDI2) with enrichment for weak discontinuities at material interfaces.
\newblock \emph{International Journal for Numerical Methods in Engineering}, 95(11):928--952, 2013.
\newblock \url{https://doi.org/10.1002/nme.4526}.

\bibitem{homel2016domaindef}
M. A. Homel, R. M. Brannon, and J. E. Guilkey.
\newblock Controlling the onset of numerical fracture in parallelized implementations of the material point method with convective particle domain interpolation domain scaling.
\newblock \emph{International Journal for Numerical Methods in Engineering}, 107(1):31--48, 2016.
\newblock \url{https://doi.org/10.1002/nme.5151}.



\bibitem{homel2016dfg}
M. A. Homel and E. B. Herbold.
\newblock Field-gradient partitioning for fracture and frictional contact in the material point method.
\newblock \emph{International Journal for Numerical Methods in Engineering}, 109(7), 2016.
\newblock \url{https://doi.org/10.1002/nme.5317}.

\bibitem{nairn2003cramp}
J. A. Nairn.
\newblock Material point method calculations with explicit cracks.
\newblock \emph{Computer Modeling in Engineering \& Sciences}, 4(6):649--663, 2003.
\newblock \url{https://doi.org/10.3970/cmes.2003.004.649}.

\bibitem{bardenhagen2000granular}
S. G. Bardenhagen, J. U. Brackbill, and D. Sulsky.
\newblock The material-point method for granular materials.
\newblock \emph{Computer Methods in Applied Mechanics and Engineering}, 187(3--4):529--541, 2000.
\newblock \url{https://doi.org/10.1016/S0045-7825(99)00338-2}.

\bibitem{bardenhagen2001contact}
S. G. Bardenhagen, J. E. Guilkey, K. M. Roessig, J. U. Brackbill, W. M. Witzel, and J. C. Foster.
\newblock An improved contact algorithm for the material point method and application to stress propagation in granular material.
\newblock \emph{Computer Modeling in Engineering \& Sciences}, 2(4):509--522, 2001.
\newblock \url{https://doi.org/10.3970/cmes.2001.002.509}.



\bibitem{nairn2020contact}
J. A. Nairn, C. C. Hammerquist, and G. D. Smith.
\newblock New material point method contact algorithms for improved accuracy, large-deformation problems, and proper null-space filtering.
\newblock \emph{Computer Methods in Applied Mechanics and Engineering}, 362:112859, 2020.
\newblock \url{https://doi.org/10.1016/j.cma.2020.112859}.

\bibitem{wallstedt2007projection}
P. C. Wallstedt and J. E. Guilkey.
\newblock Improved velocity projection for the material point method.
\newblock \emph{Computer Modeling in Engineering \& Sciences}, 19(3):223--232, 2007.

\bibitem{wallstedt2008explicit}
P. C. Wallstedt and J. E. Guilkey.
\newblock An evaluation of explicit time integration schemes for use with the generalized interpolation material point method.
\newblock \emph{Journal of Computational Physics}, 227(22):9628--9642, 2008.

\bibitem{steffen2008quadrature}
M. Steffen, R. M. Kirby, and M. Berzins.
\newblock Analysis and reduction of quadrature errors in the material point method.
\newblock \emph{International Journal for Numerical Methods in Engineering}, 76(6):922--948, 2008.
\newblock \url{https://doi.org/10.1002/nme.2360}.

\bibitem{steffen2008choices}
M. Steffen, P. C. Wallstedt, J. E. Guilkey, R. M. Kirby, and M. Berzins.
\newblock Examination and analysis of implementation choices within the material point method.
\newblock \emph{Computer Modeling in Engineering \& Sciences}, 31(2):107--128, 2008.
\newblock \url{https://doi.org/10.3970/cmes.2008.031.107}.

\bibitem{guilkey2003implicit}
J. E. Guilkey and J. A. Weiss.
\newblock Implicit time integration for the material point method: quantitative and algorithmic comparisons with the finite element method.
\newblock \emph{International Journal for Numerical Methods in Engineering}, 57:1323--1338, 2003.
\newblock \url{https://doi.org/10.1002/nme.729}.

\bibitem{hammerquist2017xpic}
C. C. Hammerquist and J. A. Nairn.
\newblock A new method for material point method particle updates that reduces noise and enhances stability.
\newblock \emph{Computer Methods in Applied Mechanics and Engineering}, 318:724--738, 2017.
\newblock \url{https://doi.org/10.1016/j.cma.2017.01.035}.

\bibitem{nairn2021fmpm}
J. A. Nairn and C. C. Hammerquist.
\newblock Material point method simulations using an approximate full mass matrix inverse.
\newblock \emph{Computer Methods in Applied Mechanics and Engineering}, 377:113667, 2021.
\newblock \url{https://doi.org/10.1016/j.cma.2021.113667}.


\bibitem{nairn2026fmpm}
J. A. Nairn.
\newblock Improved implementation of approximate full mass matrix inverse methods into material point method simulations.
\newblock arXiv:2604.07307, 2026.
\newblock \url{https://doi.org/10.48550/arXiv.2604.07307}.

\bibitem{sulsky1995solid}
D. Sulsky, S.-J. Zhou, and H. L. Schreyer.
\newblock Application of a particle-in-cell method to solid mechanics.
\newblock \emph{Computer Physics Communications}, 87(1--2):236--252, 1995.
\newblock \url{https://doi.org/10.1016/0010-4655(94)00170-7}.

\bibitem{brackbill1988flip}
J. U. Brackbill, D. B. Kothe, and H. M. Ruppel.
\newblock FLIP: a low-dissipation, particle-in-cell method for fluid flow.
\newblock \emph{Computer Physics Communications}, 48:25--38, 1988.
\newblock \url{https://doi.org/10.1016/0010-4655(88)90020-3}.

\bibitem{brackbill1992csf}
J. U. Brackbill, D. B. Kothe, and C. Zemach.
\newblock A continuum method for modeling surface tension.
\newblock \emph{Journal of Computational Physics}, 100(2):335--354, 1992.
\newblock \url{https://doi.org/10.1016/0021-9991(92)90240-Y}.


\bibitem{gan2018bspline}
Y. Gan, Z. Sun, Z. Chen, X. Zhang, and Y. Liu.
\newblock Enhancement of the material point method using B-spline basis functions.
\newblock \emph{International Journal for Numerical Methods in Engineering}, 113(3):411--431, 2018.
\newblock \url{https://doi.org/10.1002/nme.5620}.

\bibitem{leavy2019cpti}
R. B. Leavy, J. E. Guilkey, B. R. Phung, A. D. Spear, and R. M. Brannon.
\newblock A convected-particle tetrahedron interpolation technique in the material-point method for the mesoscale modeling of ceramics.
\newblock \emph{Computational Mechanics}, 64, 2019.
\newblock \url{https://doi.org/10.1007/s00466-019-01670-x}.


\bibitem{crook2025cohesive}
C. M. Crook and M. A. Homel.
\newblock A cohesive zone treatment for the material point method involving problems of large deformation and damage.
\newblock \emph{Computer Methods in Applied Mechanics and Engineering}, 448:118399, 2026.
\newblock \url{https://doi.org/10.1016/j.cma.2025.118399}.


\bibitem{xuneedleman1994}
X.-P. Xu and A. Needleman.
\newblock Numerical simulations of fast crack growth in brittle solids.
\newblock \emph{Journal of the Mechanics and Physics of Solids}, 42(9):1397--1434, 1994.
\newblock \url{https://doi.org/10.1016/0022-5096(94)90003-5}.


\bibitem{readshockley1950}
W. T. Read and W. Shockley.
\newblock Dislocation models of crystal grain boundaries.
\newblock \emph{Physical Review}, 78:275--289, 1950.
\newblock \url{https://doi.org/10.1103/PhysRev.78.275}.



\bibitem{harlow1964pic}
F. H. Harlow.
\newblock The particle-in-cell computing method for fluid dynamics.
\newblock In \emph{Methods in Computational Physics}, volume 3, pages 319--343. Academic Press, 1964.


\bibitem{weibull1951}
W. Weibull.
\newblock A statistical distribution function of wide applicability.
\newblock \emph{Journal of Applied Mechanics}, 18:293--297, 1951.

\bibitem{griffith1921}
A. A. Griffith.
\newblock The phenomena of rupture and flow in solids.
\newblock \emph{Philosophical Transactions of the Royal Society of London. Series A}, 221:163--198, 1921.

\bibitem{irwin1957}
G. R. Irwin.
\newblock Analysis of stresses and strains near the end of a crack traversing a plate.
\newblock \emph{Journal of Applied Mechanics}, 24:361--364, 1957.

\bibitem{homel2017mesoscalevalidation}
M. A. Homel, J. E. Guilkey, and R. M. Brannon.
\newblock Mesoscale validation of simplifying assumptions for modeling the plastic deformation of fluid-saturated porous material.
\newblock \emph{Journal of Dynamic Behavior of Materials}, 3:23--44, 2017.
\newblock \url{https://doi.org/10.1007/s40870-017-0092-8}.

\bibitem{homel2017porous}
M. A. Homel and E. B. Herbold.
\newblock Mesoscale modeling of porous materials using new methodology for fracture and frictional contact in the material point method.
\newblock In \emph{Dynamic Behavior of Materials, Volume 1}, Conference Proceedings of the Society for Experimental Mechanics Series, pages 97--102. Springer, 2017.
\newblock \url{https://doi.org/10.1007/978-3-319-62956-8_17}.


\bibitem{homelCrookAppleton2026sizeeffects}
M. A. Homel, C. M. Crook, and J. Appleton.
\newblock Emergent size effects in continuum damage modeling of brittle fracture with under-resolved crack-tip stresses.
\newblock Manuscript in preparation, 2026.

\bibitem{homel2022uhpc}
M. A. Homel, J. Iyer, S. J. Semnani, and E. B. Herbold.
\newblock Mesoscale model and X-ray computed micro-tomographic imaging of damage progression in ultra-high-performance concrete.
\newblock \emph{Cement and Concrete Research}, 157:106799, 2022.
\newblock \url{https://doi.org/10.1016/j.cemconres.2022.106799}.


\bibitem{malvern1969}
L. E. Malvern.
\newblock \emph{Introduction to the Mechanics of a Continuous Medium}.
\newblock Prentice-Hall, 1969.

\bibitem{ogden1997}
R. W. Ogden.
\newblock \emph{Non-Linear Elastic Deformations}.
\newblock Dover Publications, 1997.  Originally published by Ellis Horwood, 1984.

\bibitem{holzapfel2000}
G. A. Holzapfel.
\newblock \emph{Nonlinear Solid Mechanics: A Continuum Approach for Engineering}.
\newblock Wiley, 2000.

\bibitem{lekhnitskii1963}
S. G. Lekhnitskii.
\newblock \emph{Theory of Elasticity of an Anisotropic Elastic Body}.
\newblock Holden-Day, 1963.

\bibitem{simo1998inelasticity}
J. C. Simo and T. J. R. Hughes.
\newblock \emph{Computational Inelasticity}.
\newblock Springer, 1998.
\newblock \url{https://doi.org/10.1007/b98904}.

\bibitem{roache1998}
P. J. Roache.
\newblock \emph{Verification and Validation in Computational Science and Engineering}.
\newblock Hermosa Publishers, 1998.

\bibitem{boyce1988polymer}
M. C. Boyce, D. M. Parks, and A. S. Argon.
\newblock Large inelastic deformation of glassy polymers. Part I: rate dependent constitutive model.
\newblock \emph{Mechanics of Materials}, 7(1):15--33, 1988.
\newblock \url{https://doi.org/10.1016/0167-6636(88)90003-8}.

\bibitem{arruda1993boyce}
E. M. Arruda and M. C. Boyce.
\newblock A three-dimensional constitutive model for the large stretch behavior of rubber elastic materials.
\newblock \emph{Journal of the Mechanics and Physics of Solids}, 41(2):389--412, 1993.
\newblock \url{https://doi.org/10.1016/0022-5096(93)90013-6}.

\bibitem{cervera2006chiumenti}
M. Cervera and M. Chiumenti.
\newblock Mesh objective tensile cracking via a local continuum damage model and a crack tracking technique.
\newblock \emph{Computer Methods in Applied Mechanics and Engineering}, 196(1--3):304--320, 2006.
\newblock \url{https://doi.org/10.1016/j.cma.2006.04.008}.

\bibitem{borja2013plasticity}
R. I. Borja.
\newblock \emph{Plasticity: Modeling and Computation}.
\newblock Springer, 2013.
\newblock \url{https://doi.org/10.1007/978-3-642-38547-6}.



\bibitem{brannon2009kayenta}
R. M. Brannon, A. F. Fossum, and O. E. Strack.
\newblock \emph{Kayenta: Theory and User's Guide}.
\newblock Sandia National Laboratories Report SAND2009-2282, 2009.

\bibitem{homel2014arenisca}
M. A. Homel, A. Sadeghirad, D. Austin, J. Colovos, R. M. Brannon, and J. E. Guilkey.
\newblock \emph{ARENISCA: Theory and User's Guide}.
\newblock University of Utah, 2014.

\bibitem{homel2015cbtscpr}
M. A. Homel, J. E. Guilkey, and R. M. Brannon.
\newblock Numerical solution for plasticity models using consistency bisection and a transformed-space closest-point return: a nongradient solution method.
\newblock \emph{Computational Mechanics}, accepted manuscript, 2015.

\bibitem{malenda2025ghareb}
M. G. Malenda, M. A. Homel, W. M. Kibikas, C. Choens, E. Shalev, and V. Lyakhovsky.
\newblock A continuum-scale approach to predicting wellbore stability in Ghareb chalk for safe nuclear waste disposal.
\newblock In \emph{Proceedings of the 59th U.S. Rock Mechanics/Geomechanics Symposium}, ARMA 25-340, 2025.

\bibitem{brannon2017rotation}
R. M. Brannon.
\newblock \emph{Rotation, Reflection, and Frame Changes with Applications in Computational Continuum Mechanics}.
\newblock University of Utah, 2017.

\end{thebibliography}

\section{Code source inventory}
\label{sec:source-inventory}
The source inventory in Table~\ref{tab:source-inventory} lists the principal paths reviewed or extracted from the current archive.

